\section{Holonomy Return Operators}
\label{sec:holonomy-return}

Geometric interaction in NDH is governed by the connection $\nabla^{cog}$ on the
cognitive manifold $\mathcal{M}_{cog}$. When semantic excitation induces a
closed geodesic loop, parallel transport around the loop yields a holonomy
operator that encodes the cumulative geometric deformation experienced during
interaction. This section formalizes holonomy return operators and their role
in NDH interaction classification.

\subsection{Parallel Transport}

Given a geodesic $\gamma : [0,T] \to \mathcal{M}_{cog}$, parallel transport of a
tangent vector $v \in T_{\gamma(0)}\mathcal{M}_{cog}$ along $\gamma$ is defined
by the differential equation
\[
\nabla^{cog}_{\dot{\gamma}(t)} V(t) = 0,
\qquad V(0) = v.
\]
The solution $V(t)$ represents the transported vector at time $t$. Parallel
transport preserves the geometric structure induced by $\nabla^{cog}$ and
captures how interaction direction evolves along the trajectory.

\subsection{Holonomy of Closed Loops}

If $\gamma$ is a closed interaction loop satisfying $\gamma(0) = \gamma(T)$,
parallel transport defines a linear operator
\[
\mathcal{H}(\gamma) : T_x \mathcal{M}_{cog} \to T_x \mathcal{M}_{cog},
\qquad x = \gamma(0),
\]
called the \emph{holonomy return operator}. This operator represents the net
effect of transporting tangent vectors around the loop. Holonomy is trivial if
$\mathcal{H}(\gamma)$ is the identity; otherwise, the loop induces geometric
deformation.

\subsection{Modal Holonomy}

Modal-boundary crossings contribute discrete components to holonomy. If
$\gamma$ crosses a modal boundary at $t = t_0$, the tangent space undergoes a
transition governed by a boundary map
\[
B_{\gamma(t_0)} : T_{\gamma(t_0^-)}\mathcal{M}_{cog}
                 \to T_{\gamma(t_0^+)}\mathcal{M}_{cog}.
\]
The cumulative modal holonomy is the composition of all such boundary maps
encountered along the loop. Modal holonomy captures structural transitions
between recursion operators and modal adjacency effects.

\subsection{Spectral Holonomy}

Parallel transport on the spectral bundle $\mathcal{B}_S$ induces a spectral
holonomy operator
\[
\mathcal{H}_S(\gamma) : \mathcal{B}_{S,x} \to \mathcal{B}_{S,x},
\]
representing the net spectral deformation induced by traversing the loop.
Spectral holonomy arises from modal-boundary twisting, non-dual algebraic
coupling, and curvature of the interaction connection.

Spectral holonomy is nontrivial when NDH exhibits interaction-dependent spectral
drift or rotation.

\subsection{Recursion Holonomy}

The triadic recursion kernel $(R_A, R_B, R_C)$ may undergo cyclic transitions
along a closed loop. The recursion holonomy operator is defined by
\[
\mathcal{H}_R(\gamma) = R_{i_k} \circ \cdots \circ R_{i_1},
\]
where $(i_1, \dots, i_k)$ is the sequence of modal regions visited by $\gamma$.
Nontrivial recursion holonomy corresponds to interaction loops that induce
recursion-dependent structural deformation.

\subsection{Holonomy Classification}

Holonomy return operators classify interaction loops into distinct geometric
types:

\begin{itemize}
    \item \textbf{Holonomy-neutral}: $\mathcal{H}(\gamma) = \mathrm{id}$.
    \item \textbf{Modal-holonomy}: nontrivial modal boundary contribution.
    \item \textbf{Spectral-holonomy}: nontrivial spectral deformation.
    \item \textbf{Recursion-holonomy}: cyclic recursion transitions.
    \item \textbf{Mixed-holonomy}: combined modal, spectral, and recursion
    effects.
\end{itemize}

These classes form the foundation of NDH's interaction taxonomy.

\subsection{Summary}

Holonomy return operators encode the cumulative geometric deformation induced by
closed interaction loops. Modal-boundary crossings, spectral twisting, and
recursion transitions contribute to holonomy, yielding a geometric
classification of NDH interaction behavior. Holonomy is therefore a central
mechanism in NDH interaction geometry.
